% Part: proof-theory
% Chapter: natural-deduction
% Section: rules-N1

\documentclass[../../../include/open-logic-section]{subfiles}

\begin{document}

\begin{table}
\begin{center}
\begin{tabular}[t]{@{}cc@{}}
\bottomAlignProof
\AxiomC{$\Discharge{!A}{x}$}
\shortDeduce
\DeduceC{$!B$}
\DischargeRule{\Intro{\lif}}{x}
\UnaryInfC{$!A \lif !B$}
\DisplayProof
&
\bottomAlignProof
\AxiomC{$!A \lif !B$}
\AxiomC{$!A$}
\RightLabel{\Elim{\lif}}
\BinaryInfC{$!B$}
\DisplayProof
\\[4ex]
\bottomAlignProof
\AxiomC{$!A$}
\AxiomC{$!B$}
\RightLabel{\Intro{\land}}
\BinaryInfC{$!A \land !B$}
\DisplayProof
&
\begin{tabular}[b]{@{}l@{}}
\AxiomC{$!A \land !B$}
\RightLabel{\Elim{\land}[1]}
\UnaryInfC{$!A$}
\DisplayProof
\\[3ex]
\bottomAlignProof
\AxiomC{$!A \land !B$}
\RightLabel{\Elim{\land}[2]}
\UnaryInfC{$!B$}
\DisplayProof
\end{tabular}
\\[4ex]
\begin{tabular}{@{}l@{}}
\AxiomC{$!A$}
\RightLabel{\Intro{\lor}[1]}
\UnaryInfC{$!A \lor !B$}
\DisplayProof
\\[3ex]
\AxiomC{$!B$}
\RightLabel{\Intro{\lor}[2]}
\UnaryInfC{$!A \lor !B$}
\DisplayProof
\end{tabular}
&
\AxiomC{$!A \lor !B$}
\AxiomC{$\Discharge{!A}{x}$}
\shortDeduce
\DeduceC{$!C$}
\AxiomC{$\Discharge{!B}{y}$}
\shortDeduce
\DeduceC{$!C$}
\DischargeRule{\Elim{\lor}}{x\,y}
\TrinaryInfC{$!C$}
\DisplayProof
\bottomAlignProof
\\[6ex]
\bottomAlignProof
\AxiomC{$\Discharge{!A}{x}$}
\shortDeduce
\DeduceC{$\lfalse$}
\DischargeRule{\Intro{\lnot}}{x}
\UnaryInfC{$\lnot !A$}
\DisplayProof
&
\bottomAlignProof
\AxiomC{$\lnot !A$}
\AxiomC{$!A$}
\RightLabel{\Elim{\lnot}}
\BinaryInfC{$\lfalse$}
\DisplayProof
\\[4ex]
\bottomAlignProof
\AxiomC{$!A(a)$}
\RightLabel{\Intro{\lforall}}
\UnaryInfC{$\lforall[x][!A(x)]$}
\DisplayProof
&
\bottomAlignProof
\AxiomC{$\lforall[x][!A(x)]$}
\RightLabel{\Elim{\lforall}}
\UnaryInfC{$!A(t)$}
\DisplayProof
\\[4ex]
\bottomAlignProof
\AxiomC{$!A(t)$}
\RightLabel{\Intro{\lexists}}
\UnaryInfC{$\lexists[x][!A(x)]$}
\DisplayProof
&
\bottomAlignProof
\AxiomC{$\lexists[x][!A(x)]$}
\AxiomC{$\Discharge{!A(a)}{x}$}
\shortDeduce
\DeduceC{$!C$}
\DischargeRule{\Elim{\lexists}}{x}
\BinaryInfC{$!C$}
\DisplayProof
\\[4ex]
\bottomAlignProof
\AxiomC{$\lfalse$}
\RightLabel{\FalseInt}
\UnaryInfC{$!A$}
\DisplayProof
&
\bottomAlignProof
\AxiomC{$\Discharge{\lnot!A}{x}$}
\shortDeduce
\DeduceC{$\lfalse$}
\DischargeRule{\FalseCl}{x}
\UnaryInfC{$!A$}
\DisplayProof
\end{tabular}
\end{center}
\caption{\Log{N1i} અને~\Log{N1c}ના નિયમો.
$\Intro{\lforall}$માં
!!{constant}~$a$ નિષ્કર્ષમાં
અથવા કોઈ ખુલ્લી ધારણામાં
આવવો ન જોઈએ.
$\Elim{\lexists}$માં~$a$ મુખ્ય આધારવિધાન
$\lexists[x][!A(x)]$માં, નિષ્કર્ષ~$!C$માં કે
નિવૃત્ત થતી~$!A(a)$ સિવાયની કોઈ ખુલ્લી ધારણામાં
આવવો ન જોઈએ. \Log{N1i}માં
\FalseCl સામેલ નથી.}
\ollabel{tab:N1}
\end{table}
\sourcecorrection{OLREA-196}{મૂળ તકતીનો અસ્તિત્વ-નિરાકરણનો
નિયમ સ્વનિર્ધારક અચળને બધા આધારવિધાનો અને બધી
ખુલ્લી ધારણાઓમાંથી બાકાત રાખે છે. પરંતુ ગૌણ
ઉપસાબિતીની નિવૃત્ત થતી ધારણા $!A(a)$માં તે આવવો
જ પડે છે. મુખ્ય આધારવિધાન, નિષ્કર્ષ અને બાકી
ખુલ્લી ધારણાઓ અંગેની યોગ્ય શરત મૂકી છે; તે
મુખ્ય પાઠની સ્વનિર્ધારક-ચલની શરત સાથે પણ મેળ ખાય છે.}

\end{document}
